\section{Discussion}

\subsection{Training and Evaluation}
The leakage-aware model trained smoothly, with training and validation accuracy
rising together and no sign of overfitting. It is worth being clear about how we
read this. In our first run, validation accuracy jumped above 99\% within a few
epochs and tracked training accuracy almost exactly, and we first took the tiny gap
as proof of strong generalization. The cross-dataset analysis showed that reading
was wrong. Under a leaky random split, a small gap is actually expected, because
near-duplicate frames sit on both sides of the split, so the validation set is not
really new to the model. Only after switching to the leakage-aware split does a
small gap mean what we want it to mean, and there the model reaches 94.33\% on a
genuinely unseen test set. The label smoothing we used ($\epsilon = 0.10$) also
keeps the loss from reaching zero, which is why the loss settles near $0.37$ even at
high accuracy.

\subsection{Per-Class Performance}
Accuracy is high and even across the alphabet: every class reaches an F1-score above
0.85, and none collapses. The weaker classes are \texttt{kaaf}, \texttt{dal}, and
\texttt{thal}, with \texttt{fa} and \texttt{ra} slightly behind the rest. These are
the letters whose hand shapes look almost the same and differ only in small details.
A clear example is \texttt{dal} and \texttt{thal}, which differ only by a dot. Such
cases are hard for any system that works from a single still image, because the
distinguishing detail is small and not always visible.

\subsection{Misclassification Analysis}
Looking at the wrong predictions, most errors happen between letters with very
similar hand shapes rather than between random pairs, which matches the per-class
results above. The model is also less confident on these cases, with lower
probability scores than on correct predictions. That suggests a simple improvement
for future versions: when confidence is low, the system could ask the user to repeat
the sign instead of guessing.

One earlier case is worth noting. In the cross-dataset test, the pair
\texttt{ta}/\texttt{taa} failed completely, which traced back to a difference in how
the two datasets label or sign that pair, not to a weakness in the model. After
retraining within a single, consistent dataset, both letters are recognized well. It
is a good reminder that matching letters by name across datasets is not the same as
matching the actual gesture.

\subsection{Strengths of the Approach}
A few choices helped the system work well:
\begin{itemize}
  \item Starting from ImageNet pretrained weights gave a strong base, so the model
        adapted to hand images with limited data.
  \item The two-phase training kept the pretrained features intact early on and let
        the network adjust gradually.
  \item MobileNetV2 is light enough to run in real time on a normal laptop, with no
        special hardware.
  \item The spelling tutor turns the model into something a learner can practice
        with, not just a classifier.
  \item Most importantly, we evaluated the model honestly. Instead of reporting a
        single inflated benchmark number, we found the leakage, measured the model on
        a separate dataset, and retrained so that the reported 94.33\% reflects real
        performance.
\end{itemize}

\subsection{Limitations}
Several limitations remain:
\begin{itemize}
  \item \textbf{Static letters only.} The system recognizes single letter signs, not
        dynamic word-level signs that involve motion. Those would need a temporal
        model.
  \item \textbf{Similar letter pairs.} Letters that differ only by a small detail stay
        the main source of errors and may need extra cues to separate.
  \item \textbf{Still a domain gap.} Even after retraining, accuracy across a fully
        different dataset is lower than within one dataset. Real users will vary in
        ways no single dataset fully covers.
  \item \textbf{One class dropped.} The color dataset has no \texttt{yaa}, so the
        retrained model covers 31 of the 32 letters. Adding it back needs more data.
  \item \textbf{Single-frame decisions.} Each frame is classified on its own, which
        can make predictions jump while the user moves between signs. Averaging over a
        few frames helps but does not fully fix it.
\end{itemize}

\subsection{Broader Impact}
This work targets a group with few tools available to them: Arabic-speaking deaf and
hard-of-hearing learners. By giving immediate feedback during letter-spelling
practice, ArSignTutor makes self-directed learning easier, and it can help early
learners, parents of deaf children, and teachers in places without specialized
instructors. Just as important, we hope the leakage analysis is a useful reminder for
similar projects: a near-perfect score is not always a good model, and it is worth
checking how the data was split before trusting the number.
